% ======================================================================
% Closed-Form Fisher Information for Homogeneous Weibull Series Systems
% with Masked Failure Causes
%
% Target: Lifetime Data Analysis
% ======================================================================
\documentclass[12pt]{article}

% --- Packages ---
\usepackage[margin=1in]{geometry}
\usepackage{amsmath,amssymb,amsthm}
\usepackage{bm}
\usepackage{natbib}
\usepackage{booktabs}
\usepackage[colorlinks,citecolor=blue,linkcolor=blue,urlcolor=blue]{hyperref}

% --- Theorem environments ---
\newtheorem{theorem}{Theorem}
\newtheorem{corollary}[theorem]{Corollary}
\theoremstyle{remark}
\newtheorem{remark}{Remark}

% --- Notation ---
\newcommand{\lsys}{\lambda_{\mathrm{sys}}}
\newcommand{\bsys}{\beta_{\mathrm{sys}}}
\newcommand{\ind}{\mathbf{1}}
\newcommand{\ee}{\mathbf{e}}

\title{Closed-Form Fisher Information for Homogeneous Weibull Series
  Systems with Masked Failure Causes}

\author{Alexander R.\ Towell\\
  \small Department of Mathematics and Statistics\\
  \small Southern Illinois University Edwardsville\\
  \small \texttt{lex@metafunctor.com}\\
  \small ORCID: 0000-0001-6443-9897}

\date{}

\begin{document}
\maketitle

\begin{abstract}
We derive the expected per-observation Fisher information matrix in
closed form for a series system of $m$ independent Weibull components
sharing a common shape parameter.  In the rate parametrization
$\lambda_j = \beta_j^{-k}$, the matrix entries involve only the
Euler--Mascheroni constant~$\gamma$ and~$\pi^2/6$; no numerical
integration is required.  We prove that masking---observing a candidate
set of possible failure causes rather than the true cause---leaves all
shape-related Fisher information unchanged.  The entire information loss
from masking is confined to the rate block of the Fisher information
matrix, where it equals the information loss derived for the exponential
special case.  Monte Carlo simulation with $10^6$ replications confirms
all analytical entries to within 0.5\% relative error.
\end{abstract}

\noindent\textbf{Keywords:}
Fisher information $\cdot$
Weibull distribution $\cdot$
series system $\cdot$
competing risks $\cdot$
masked failure cause


% ======================================================================
\section{Introduction}\label{sec:intro}
% ======================================================================

Series systems, in which the system fails when any component fails,
arise naturally in reliability engineering.  When components have
independent lifetimes, the system failure time is the minimum of the
component lifetimes, placing this problem within the classical competing
risks framework~\citep{crowder2001}.  A common complication in practice
is \emph{masking}: the failed component is not directly observed, and
only a candidate set of possible causes is
reported~\citep{miyakawa1984}.

\citet{towell2025masked} developed a likelihood framework for masked
series systems under three regularity conditions (C1--C3) that unify
right-censoring, masking, and their interaction.
\citet{towell2026exponential} specialized this framework to exponential
components and derived the Fisher information matrix in closed form,
establishing a baseline for studying information loss under masking.

The Weibull distribution is the workhorse of reliability
modeling~\citep{lawless2003, meeker1998}.  When all components share a
common shape parameter---the \emph{homogeneous} Weibull model---the
system lifetime is itself Weibull and the failed component is independent
of the failure time, yielding a tractable
likelihood~\citep{pareek2009}.  Existing treatments derive score
equations and compute the \emph{observed} Fisher information numerically
at the maximum likelihood estimate~\citep{pareek2009, kundu2003}.  To
our knowledge, the \emph{expected} Fisher information matrix has not been
obtained in closed form for this model.

This paper derives the expected per-observation Fisher information matrix
for the homogeneous Weibull series system in the rate parametrization
$\lambda_j = \beta_j^{-k}$.  Our contributions are:
%
\begin{enumerate}
\item \textbf{Closed-form FIM} (Theorem~\ref{thm:fim}): All
  $(m{+}1)^2$ entries of the Fisher information matrix reduce to
  elementary functions of the Euler--Mascheroni constant~$\gamma$ and
  $\pi^2/6$.  No numerical integration is required.

\item \textbf{Masking invariance} (Theorem~\ref{thm:masking}): Under
  any masking mechanism satisfying the C2 symmetry condition of
  \citet{towell2025masked}, the shape information~$I_{kk}$ and the
  shape--rate cross-information~$I_{k,\lambda_j}$ are unchanged.  Only
  the rate--rate block is affected.

\item \textbf{Structural decomposition}: The masked Fisher information
  matrix equals the unmasked matrix with its rate block replaced by the
  exponential masked-data Fisher information
  of~\citet{towell2026exponential}, evaluated at rates
  $\lambda_j = \beta_j^{-k}$.
\end{enumerate}

Section~\ref{sec:model} defines the model.
Section~\ref{sec:fim} derives the Fisher information matrix.
Section~\ref{sec:masking} establishes masking invariance.
Section~\ref{sec:verification} reports numerical verification.
Section~\ref{sec:discussion} concludes.


% ======================================================================
\section{Model}\label{sec:model}
% ======================================================================

Consider a series system of $m$ independent components.  Component~$j$
has lifetime $T_j \sim \mathrm{Weibull}(k, \beta_j)$ with common shape
$k > 0$ and component-specific scale $\beta_j > 0$.  The system fails at
time $T = \min(T_1, \ldots, T_m)$, and the failed component is
$K = \arg\min_j T_j$.

\subsection{Rate parametrization}

Define the rate parameters $\lambda_j = \beta_j^{-k}$ and the system
quantities
%
\begin{equation}\label{eq:system}
\lsys = \sum_{j=1}^m \lambda_j, \qquad
\bsys = \lsys^{-1/k}, \qquad
w_j = \frac{\lambda_j}{\lsys}.
\end{equation}
%
The system hazard function is $h(t) = k\,\lsys\,t^{k-1}$, additive in
the~$\lambda_j$.  The cause probabilities~$w_j$ do not depend on~$k$,
which is the key to the masking invariance of
Theorem~\ref{thm:masking}.

\subsection{Joint density and independence}

The joint density of the system failure time~$T$ and failed
component~$K$ factorizes as
%
\begin{equation}\label{eq:joint}
f(t, K{=}j \mid k, \bm{\lambda})
= \underbrace{k\,\lsys\,t^{k-1}
  \exp\!\bigl(-\lsys\,t^k\bigr)}_{f_T(t)}
  \;\cdot\; \underbrace{w_j\vphantom{t^{k-1}}}_{P(K=j)}.
\end{equation}
%
Thus $T$ and $K$ are independent, with
$T \sim \mathrm{Weibull}(k, \bsys)$ and
$K \sim \mathrm{Categorical}(w_1, \ldots, w_m)$.  This factorization,
a direct consequence of the common shape, is well known in the competing
risks literature~\citep{crowder2001, pareek2009}.

\subsection{Log-likelihood}

The per-observation log-likelihood for an exact, unmasked observation
$(t, K{=}j)$ is
%
\begin{equation}\label{eq:loglik}
\ell_1(k, \bm{\lambda})
= \log k + \log \lambda_j + (k-1)\log t - \lsys\,t^k.
\end{equation}


% ======================================================================
\section{Fisher information}\label{sec:fim}
% ======================================================================

\subsection{Score equations}

Differentiating~\eqref{eq:loglik} with respect to the parameters
$\theta = (k, \lambda_1, \ldots, \lambda_m)$ gives
%
\begin{align}
s_k &= \frac{1}{k} + (1 - V)\log t, \label{eq:score_k}\\
s_{\lambda_\ell} &= \frac{\ind(K = \ell)}{\lambda_\ell} - t^k,
  \label{eq:score_lam}
\end{align}
%
where $V = \lsys\,t^k$.  Since
$T \sim \mathrm{Weibull}(k, \bsys)$, the transformation
$V = \lsys\,T^k$ yields $V \sim \mathrm{Exp}(1)$, and
%
\begin{equation}\label{eq:logT}
\log T = \log \bsys + \frac{\log V}{k}.
\end{equation}

\subsection{Closed-form Fisher information matrix}

\begin{theorem}\label{thm:fim}
Let\/ $\gamma \approx 0.5772$ denote the Euler--Mascheroni constant and
define the centering constant $\mu = \log \bsys + (1-\gamma)/k$.  The
per-observation expected Fisher information matrix for the homogeneous
Weibull series system in the parametrization
$\theta = (k, \lambda_1, \ldots, \lambda_m)$ is
%
\begin{equation}\label{eq:fim}
I(\theta)
= \begin{pmatrix}
  \mu^2 + \dfrac{\pi^2}{6k^2}
    & \dfrac{\mu}{\lsys}\,\ee^\top \\[10pt]
  \dfrac{\mu}{\lsys}\,\ee
    & D
  \end{pmatrix},
\end{equation}
%
where $\ee = (1, \ldots, 1)^\top \in \mathbb{R}^m$ and
$D = \operatorname{diag}\!\bigl(
  (\lambda_1\lsys)^{-1}, \ldots, (\lambda_m\lsys)^{-1}\bigr)$.
\end{theorem}

\begin{proof}
The second derivatives of~\eqref{eq:loglik} are
%
\begin{align}
\frac{\partial^2 \ell_1}{\partial k^2}
  &= -\frac{1}{k^2} - V(\log t)^2,
  \label{eq:hess_kk} \\
\frac{\partial^2 \ell_1}{\partial \lambda_\ell^2}
  &= -\frac{\ind(K = \ell)}{\lambda_\ell^2},
  \label{eq:hess_ll} \\
\frac{\partial^2 \ell_1}{\partial \lambda_\ell \,\partial \lambda_r}
  &= 0 \quad (\ell \neq r),
  \label{eq:hess_lr} \\
\frac{\partial^2 \ell_1}{\partial k\,\partial \lambda_\ell}
  &= -t^k \log t.
  \label{eq:hess_kl}
\end{align}
%
We compute $I(\theta) = -E[\nabla^2 \ell_1]$ using the substitution
$V = \lsys\,T^k \sim \mathrm{Exp}(1)$ and the standard moments
%
\begin{equation}\label{eq:moments}
E[V \log V] = 1 - \gamma, \qquad
E[V(\log V)^2] = (1-\gamma)^2 + \frac{\pi^2}{6} - 1.
\end{equation}

\textit{Rate block.}\enspace
Since $P(K = \ell) = w_\ell = \lambda_\ell/\lsys$,
\[
I_{\lambda_\ell,\lambda_\ell}
= \frac{w_\ell}{\lambda_\ell^2}
= \frac{1}{\lambda_\ell\,\lsys},
\]
and $I_{\lambda_\ell,\lambda_r} = 0$ for $\ell \neq r$.

\textit{Cross terms.}\enspace
Using $T^k = V/\lsys$ and~\eqref{eq:logT},
\[
I_{k,\lambda_\ell}
= E[T^k \log T]
= \frac{1}{\lsys}\Bigl[
  \log \bsys + \frac{1-\gamma}{k}
\Bigr]
= \frac{\mu}{\lsys}.
\]

\textit{Shape entry.}\enspace
From~\eqref{eq:hess_kk},
$I_{kk} = k^{-2} + E[V(\log T)^2]$.
Expanding $(\log T)^2$ via~\eqref{eq:logT} and
applying~\eqref{eq:moments}:
%
\begin{align}
I_{kk}
&= \frac{1}{k^2}
 + (\log \bsys)^2
 + \frac{2(1-\gamma)}{k}\log \bsys
 + \frac{(1-\gamma)^2 + \pi^2/6 - 1}{k^2}
\notag\\
&= \Bigl(\log \bsys + \frac{1-\gamma}{k}\Bigr)^{\!2}
 + \frac{\pi^2}{6k^2}
= \mu^2 + \frac{\pi^2}{6k^2}.
\tag*{\qedhere}
\end{align}
\end{proof}

\begin{remark}[Special cases]\label{rem:special}
\begin{enumerate}
\item[\textup{(i)}]
For $m = 1$, transforming to the $(k, \beta)$ parametrization via
$\lambda = \beta^{-k}$ recovers the classical single-Weibull Fisher
information
$I_{kk}^{(k,\beta)} = [(1-\gamma)^2 + \pi^2/6]\,/\,k^2$
\citep{lawless2003}.

\item[\textup{(ii)}]
For $k = 1$ (exponential components), $\lambda_j = 1/\beta_j$ and the
rate block~$D$ reduces to the unmasked exponential Fisher information
of~\citet{towell2026exponential}.

\item[\textup{(iii)}]
The rate block is diagonal: information about distinct component rates is
orthogonal.  Each rate~$\lambda_j$ contributes information only through
the indicator $\ind(K{=}j)$, which identifies component~$j$ as the
failure cause.

\item[\textup{(iv)}]
The cross terms $I_{k,\lambda_\ell} = \mu/\lsys$ are identical for
all~$\ell$: the shape parameter couples symmetrically to all component
rates through the system failure time.
\end{enumerate}
\end{remark}


% ======================================================================
\section{Masking invariance}\label{sec:masking}
% ======================================================================

Under the C2 masking condition of~\citet{towell2025masked}, the observer
reports a candidate set $C \subseteq \{1, \ldots, m\}$ containing the
true failed component~$K$, rather than~$K$ itself.  The conditional
probability $P(C \mid K{=}j)$ depends on the masking mechanism but not
on the model parameters $(k, \bm{\lambda})$.

The per-observation log-likelihood for a masked observation $(t, C)$ is
%
\begin{equation}\label{eq:loglik_masked}
\ell_1^{\mathrm{masked}}
= \log k + (k-1)\log t - \lsys\,t^k
  + \log\!\Bigl(\sum_{j \in C} \lambda_j\,P(C \mid K{=}j)\Bigr).
\end{equation}

\begin{theorem}\label{thm:masking}
Under any masking mechanism satisfying~C2, the shape information
$I_{kk}$ and the shape--rate cross-information $I_{k,\lambda_\ell}$ are
identical to their unmasked values.  The masked Fisher information matrix
has the block structure
%
\begin{equation}\label{eq:fim_masked}
I^{\mathrm{masked}}(\theta)
= \begin{pmatrix}
  \mu^2 + \dfrac{\pi^2}{6k^2}
    & \dfrac{\mu}{\lsys}\,\ee^\top \\[10pt]
  \dfrac{\mu}{\lsys}\,\ee
    & I^{\mathrm{exp}}_{\lambda\lambda}
  \end{pmatrix},
\end{equation}
%
where $I^{\mathrm{exp}}_{\lambda\lambda}$ is the rate block of the
exponential masked-data Fisher information
from~\citet{towell2026exponential}, evaluated at rates
$\lambda_j = \beta_j^{-k}$.
\end{theorem}

\begin{proof}
The final term in~\eqref{eq:loglik_masked} depends on~$\bm{\lambda}$
but not on~$k$.  Therefore
\[
\frac{\partial^2 \ell_1^{\mathrm{masked}}}{\partial k^2}
= -\frac{1}{k^2} - V(\log t)^2,
\qquad
\frac{\partial^2 \ell_1^{\mathrm{masked}}}{\partial k\,\partial
  \lambda_\ell}
= -t^k \log t,
\]
which are identical to the unmasked
case~\eqref{eq:hess_kk}--\eqref{eq:hess_kl}.  Taking expectations
yields the invariance of $I_{kk}$ and $I_{k,\lambda_\ell}$.

For the rate block, the Hessian of
$-\lsys\,t^k$ with respect to $\lambda_\ell,\lambda_r$ is zero;
the rate--rate information comes entirely from the
$\log(\cdot)$ term in~\eqref{eq:loglik_masked}, which involves
only the rates~$\bm{\lambda}$ and the candidate set~$C$.  Since
$K \perp T$ in the homogeneous model, $C$~is independent of~$T$, and
the distribution of~$C$ depends on~$\bm{\lambda}$ only through the
cause probabilities $w_j = \lambda_j/\lsys$---exactly as in the
exponential case.
\end{proof}

\begin{remark}[Information loss]\label{rem:info_loss}
The difference
\[
I^{\mathrm{complete}} - I^{\mathrm{masked}}
= \begin{pmatrix}
  0 & \bm{0}^\top \\
  \bm{0} & D - I^{\mathrm{exp}}_{\lambda\lambda}
  \end{pmatrix}
\]
is zero everywhere except in the rate block, and positive semidefinite by
the data processing inequality.  Masking obscures component identity but
preserves the failure time, degrading information about the individual
rates while leaving shape information intact.
\end{remark}

\begin{remark}[Parametrization dependence]\label{rem:param}
The invariance of Theorem~\ref{thm:masking} is specific to the rate
parametrization $\lambda_j = \beta_j^{-k}$.  In the
$(k, \beta_1, \ldots, \beta_m)$ parametrization, the cause
probabilities $w_j = \beta_j^{-k}/\sum_\ell \beta_\ell^{-k}$ depend
on~$k$, so the masking term in~\eqref{eq:loglik_masked} couples to~$k$
and the invariance breaks down.  The rate parametrization is natural
precisely because it decouples component identity from the time scale.
\end{remark}


% ======================================================================
\section{Numerical verification}\label{sec:verification}
% ======================================================================

We verify Theorem~\ref{thm:fim} by Monte Carlo simulation
in~R~\citep{r2025}.  For each replication, we generate a system failure
time~$T$ and failed component~$K$ from the joint
density~\eqref{eq:joint}, compute the score
vector~\eqref{eq:score_k}--\eqref{eq:score_lam} and
Hessian~\eqref{eq:hess_kk}--\eqref{eq:hess_kl}, and accumulate
two estimators: the outer product of scores and the negative Hessian.

\subsection{Primary test case}

Table~\ref{tab:mc} reports results for $m = 3$ components with
$k = 1.5$ and $\bm{\beta} = (100, 200, 150)$, using $n = 10^6$
replications.  The derived quantities are
$\bm{\lambda} = (0.001, 0.0000354, 0.000544)$,
$\lsys = 0.001578$, $\bsys = 66.51$, and $\mu = 4.47$.

\begin{table}[ht]
\centering
\caption{Analytical FIM entries (Theorem~\ref{thm:fim}) versus Monte
  Carlo estimates ($n = 10^6$).  Relative error is reported for nonzero
  entries; the off-diagonal rate entries $I_{\lambda_\ell,\lambda_r}$
  ($\ell \neq r$) are analytically zero and confirmed numerically.}
\label{tab:mc}
\smallskip
\begin{tabular}{@{}lcccc@{}}
\toprule
Entry & Analytical & MC (scores) & MC (Hessian) & Max rel.\ error \\
\midrule
$I_{kk}$                & 20.22 & 20.20 & 20.22 & 0.10\% \\
$I_{k,\lambda_\ell}$    & 2833  & 2832  & 2834  & 0.04\% \\
$I_{\lambda_1,\lambda_1}$ & 634.0 & 633.4 & 633.7 & 0.09\% \\
$I_{\lambda_2,\lambda_2}$ & 17\,900 & 17\,890 & 17\,900 & 0.06\% \\
$I_{\lambda_3,\lambda_3}$ & 1165  & 1164  & 1164  & 0.09\% \\
\bottomrule
\end{tabular}
\end{table}

All nonzero entries agree to within 0.5\% relative error across both
estimators.

\subsection{Special cases}

\textit{Single component ($m = 1$).}\enspace
With $k = 2$ and $\beta = 50$, the Jacobian-transformed
$I_{kk}^{(k,\beta)}$ from Theorem~\ref{thm:fim} matches the standard
Weibull result $[(1-\gamma)^2 + \pi^2/6]/k^2$ to machine precision.

\textit{Exponential ($k = 1$).}\enspace
With $\bm{\beta} = (100, 200, 150)$, the rate block of
Theorem~\ref{thm:fim} matches the unmasked exponential FIM
$I_{\lambda_\ell,\lambda_\ell} = 1/(\lambda_\ell\,\lsys)$ to machine
precision.


% ======================================================================
\section{Discussion}\label{sec:discussion}
% ======================================================================

We derived the expected Fisher information matrix in closed form for a
homogeneous Weibull series system with masked failure causes.  The result
(Theorem~\ref{thm:fim}) involves only the Euler--Mascheroni constant
and~$\pi^2/6$, extending the known single-Weibull FIM to the
multi-component series setting.  The masking invariance
(Theorem~\ref{thm:masking}) shows that masking degrades only the
component identification information---the rate block---while leaving
shape information untouched.

\paragraph{Applications.}
The closed-form FIM enables Cram\'er--Rao lower bound calculations and
$D$-optimal experimental design for Weibull series systems without
numerical integration.  For mixed observation types (exact, masked,
right-censored), the per-observation contributions differ in form but the
exact-data FIM provides an upper bound on the information per
observation.

\paragraph{Extensions.}
For right-censored observations at time~$c$, the score contributions
$s_k^{\mathrm{RC}} = -\lsys\,c^k \log c$ and
$s_{\lambda_\ell}^{\mathrm{RC}} = -c^k$ are available in closed form,
but the expected FIM requires specification of the censoring
distribution~\citep{escobar1994}.  For the heterogeneous Weibull model
(distinct shapes $k_1, \ldots, k_m$), the independence $K \perp T$
fails, and closed-form expressions are unlikely; numerical computation of
the observed FIM remains the practical
approach~\citep{pareek2009}.

\paragraph{Software.}
The analytical FIM of Theorem~\ref{thm:fim} and the Monte Carlo
verification of Section~\ref{sec:verification} are implemented in~R.
Source code is available from the author.


% ======================================================================
% References
% ======================================================================
\bibliographystyle{plainnat}
\bibliography{refs}

\end{document}
